\section{Paying shared raw support once}\label{sec:shared-support}

The preceding covers estimate source gap area without adding revelations.
Here we use a different mechanism: reveal a raw source set shared by many
outputs, pay its union once, and apply the one-exception area bound to
the sources that remain. All assertions in this section use $q=2$.

\subsection{A finite global-sharing bound}

Retain the notation $A_\rho$, $R_{e,\rho}$ and $Z_\rho(t)$ of
\cref{eq:older-closure,eq:two-exception-decomposition}. Fix
$J\ge R\ge\max(K,1)$, $M=2R+2$, and set
\begin{equation}\label{eq:sharing-constants}
 \ell=4R+7,\quad C_R=C_0(1+3C_0)^2,\quad
 B_R=\ell(\ell+1),\quad A_R=\beta NC_Rc,
 \qquad C_0=16+24M(M-1).
\end{equation}
Let $\mathcal T_{\rm out}$ be the set of all $L$ actual output-bit times,
and define
\[
 Z_\rho^{\rm all}=\bigcup_{t\in\mathcal T_{\rm out}}Z_\rho(t),\qquad
 G_\rho=\left|\max_{\succeq_e}Z_\rho^{\rm all}\right|,\qquad
 \overline G_{J,R}=\mathbb E_\rho G_\rho.
\]
An empty seed set has $G_\rho=0$. The order $\succeq_e$ is complete
older reachability, not merely temporal order. These maximal seeds are
geometric witnesses, not independent records of bounded bit capacity.
Let
\begin{equation}\label{eq:global-newer-support}
 S^{\rm new}_\rho
 =x\left(\bigcup_{z\in Z_\rho^{\rm all}}R_{e,\rho}(z)\right).
\end{equation}
Transitivity and maximality cover this entire set by the older closures
of the $G_\rho$ maximal seeds. Each has at most $B_R$ endpoints, so
$|S^{\rm new}_\rho|\le B_RG_\rho$.

\begin{proposition}[The full shared-support payment]\label{prop:finite-sharing}
\formal{SourceCovers/GlobalSharing.lean}{SourceCovers.shared_support_count_deficit,SourceCovers.shared_support_paid_deficit,SourceCovers.phase_shared_deficit,SourceCovers.sharedRawSupport,SourceCovers.sharedSeedCount,SourceCovers.exactSourceBank_shared_cover,SourceCovers.sharedRawSupport_card_le_sum,SourceCovers.sharedRawSupport_blocks_card_le,SourceCovers.exists_paidBank_shared_deficit}
With $\alpha,\tau,a$ as in \cref{sec:source-covers}, for every integer
$w\ge1$,
\begin{equation}\label{eq:finite-sharing}
 \frac\Delta L\ge
 \frac{(1-\alpha)_+^2}{2(2K+1)(\ell+1)}
 -\frac{wB_R\overline G_{J,R}}L
 -\mathsf E_2(w;J,2R+2,R,a,\tau).
\end{equation}
The additional raw revelation has at most $wB_RG_\rho$ bits for every
grid shift, so the residual-cut shift may be chosen using the original
complete two-exception bank.

More sharply, let $\overline S=\mathbb E_\rho|S^{\rm new}_\rho|$ and
$C_1=C_0(1+3C_0)$. Define $\mathsf E_{1\mid2}$ by replacing only the
first movement term $2C_R(T-1)/(dL)$ in $\mathsf E_2$ by
$2C_1(T-1)/(dL)$; its residual-cut term still uses $C_R$. Then
\[
 \frac\Delta L\ge
 \frac{(1-\alpha)_+^2}{2(2K+1)(\ell+1)}
 -\frac{w\overline S}L-\mathsf E_{1\mid2}.
\]
Neither finite bound requires a linear-runtime hypothesis. The formal
statements also allow any $M\ge2$ with $R<M$, taking $\ell=2M+3$;
the distinguished-family budget $K$ and protection budget $R$ may be
independent.
\end{proposition}
\begin{outline}
Fix a phase and select its residual shift $\theta$ for the complete
bank. In fixed raw-coordinate order reveal every bit in
\[
 U_{\rho,\theta}=
       \bigcup_{y\in S^{\rm new}_\rho}B_\theta(y).
\]
For every $\theta$, the finite-union and same-block cardinality bounds give
\[
 |U_{\rho,\theta}|\le\min\{L,w|S^{\rm new}_\rho|\}\le wB_RG_\rho.
\]
Complete older closure is unchanged when its seeds
are replaced by maximal seeds, so the raw-support count also retains
the sum of their actual closure sizes before using $B_RG_\rho$.
The source-replacement lemma applies whenever the full bank's sources
are contained in the remaining bank's sources and a common raw set.
Equal original transcripts, equal revealed raw bits, and equal remaining
source blocks then imply equal full-bank source blocks. Applying the original
recovery theorem therefore gives every output from the older bank and
one common joint transcript on the same correctness set. Its frozen
source gap has the one-exception lower bound
$(1-\alpha)_+^2/[2(2K+1)(\ell+1)]$ by
\cref{lem:fractional-gap}. Use the older bank for these sources;
its summed variation is at most $C_1(T-1)$, while the full-bank
variation in the residual-cut payment is at most $C_R(T-1)$.
\Cref{thm:cyclic-freezing} applies with total information charge
$\beta|\Gamma|+|U_{\rho,\theta}|$, without independence assumptions.
The intermediate formal bound keeps the actual union cardinality and
both actual variations. Average the complete signed inequality over
the same phases and use the existing average rank-payment bound.
This gives the sharper statement; $\overline S\le B_R\overline G_{J,R}$
and $C_1\le C_R$ give \cref{eq:finite-sharing} directly.
Revealing a seed vertex itself
would not reveal its entire source closure; the raw bits above are
what have actually been paid for.
\end{outline}

This supplies a verified finite certificate when many outputs share
complete raw support. It does not assert that such support is small
in an arbitrary tape computation; the following criteria isolate
the quantitative bounds that would suffice.

\subsection{Uniform sharing criteria}

For the uniform statements below, select the first $J$ original core
families in the hop ordering of \cref{sec:uniform}, protect the first
$R$, and rebuild the selected core's forest and ranks before choosing
a phase. All quantities indexed by $J,R$, including $\overline G_{J,R}$
and the raw covering numbers below, refer to these prefix selections.
Absent families are omitted.

Let $(\eta_J)_{J\ge1}$ be a fixed positive sequence tending to zero that
bounds the normalized size of the first omitted original core family:
$|G_{J+1}|/L\le\eta_J$, uniformly over the class. The universal choice
$\eta_J=c/(J+1)$ is available from sorted hop counts. Whenever another
bound is given, it may be replaced by its minimum with $c/(J+1)$.

\begin{theorem}[Theorem H: uniform exclusion by global sharing]\label{thm:uniform-sharing}
\formal{SourceCovers/UniformSharing.lean}{SourceCovers.sharing_uniform_deficit,SourceCovers.liminf_sharing_uniform_deficit,SourceCovers.sublinear_sharing_uniform_deficit,SourceCovers.power_sharing_uniform_deficit,SourceCovers.AdmissibleGeometry.HasSharingBounds,SourceCovers.NearFullSequence.JointProfile.prefix_sharing_limit,SourceCovers.NearFullSequence.JointProfile.sharing_tail_lower}
Fix $c,\beta,N$, a common sequence $\eta_J$ as above, and an unbounded
set $\mathcal R\subseteq\{1,2,\ldots\}$. For each $R\in\mathcal R$
fix a finite nonnegative function $g_R(J)$ on $J\ge R$ such that
\begin{equation}\label{eq:sharing-smallness}
 \liminf_{J\to\infty}\eta_Jg_R(J)=0.
\end{equation}
Consider admissible two-exception geometries for which
$\overline G_{J,R}\le g_R(J)$ for every $R\in\mathcal R$ and $J\ge R$.
These bounds and index sets are fixed independently of dimension.
There exist $\eps>0,d_0$ such that every such geometry with
$d\ge d_0$ and $a_0\le\eps L$ satisfies $\Delta\ge\eps L$, equivalently
$|\mathcal C|\le2^{(1-\eps)L}$.
The same conclusion holds for nonnegative $\eta_J$ tending to zero.

In particular it suffices, without a size envelope, that
$\liminf_J g_R(J)/(J+1)=0$ for each $R\in\mathcal R$.
Under $|G_i|\le HL/i^p$ with fixed $H\ge0$ and $p>0$, it suffices that
$\liminf_J g_R(J)/(J+1)^p=0$; this includes subquadratic sharing
for the inverse-square envelope.
\end{theorem}
\begin{outline}
Negating the conclusion gives a near-full sequence with a joint profile.
Choose a finite $K$ with $V_K<1$ and fix $R\in\mathcal R$.
For each finite $J\ge R$ use
$a=\max(a_0,\lceil\eta_J L\rceil)$ and
$w=\lceil\lambda L\rceil$, with $\lambda>0$ fixed before dimension.
Take $J\ge K$ as well. The finite sharing bound permits independent
$K,R$ and any positive integer width. Its dimension limit gives
\begin{equation}\label{eq:sharing-prefix-limit}
 \frac{(1-V_{K,J})^2}{2(2K+1)(\ell+1)}
 -\frac{\beta\tau_{R,J}}{R+2}
 \le \lambda[2+B_Rg_R(J)]+\frac{A_R\eta_J}\lambda.
\end{equation}
The finite sums $V_{K,J},\tau_{R,J}$ are as in
\cref{eq:finite-prefix-limit}, and $V_{K,J}\le V_K<1$.
Let $\mu_{J,R}$ be the left side and put $D_J=3+B_Rg_R(J)>0$.
Increase the coefficient $2$ to $3$. Whenever $\mu_{J,R}>0$, test at
\[
 \lambda=\frac{\mu_{J,R}}{2D_J}
 \quad\text{to obtain}\quad
 \mu_{J,R}^2\le4A_R\eta_JD_J.
\]
This uses no positive lower bound on $A_R$ or $\eta_J$ and no square-root
minimizer. The right side tends to zero along arbitrarily late cutoffs
provided by \cref{eq:sharing-smallness}, since $\eta_J\to0$.
The finite sums converge along those cutoffs, so their limiting
credit cannot be positive. As $\ell+1=4(R+2)$, this gives
\[
 \beta\tau_R\ge\frac{(1-V_K)^2}{8(2K+1)}.
\]
This holds for every $R\in\mathcal R$, contradicting
$\tau_R\to0$. The two special cases use respectively
$\eta_J=c/(J+1)$ and $\eta_J=H/(J+1)^p$.
Any $p>0$ makes the latter tend to zero; summability of the actual
profile was already proved and need not follow from the envelope.
The formal lower limits take values in $[0,\infty]$; zero is equivalent
to arbitrarily late values below every positive threshold.
\end{outline}

The bound $g_R$ may grow arbitrarily fast with the fixed protected
scale $R$. Its uniformity in dimension is essential: for a single
finite geometry, increasing $J$ beyond all its families proves no
uniform statement. A common unbounded set of cutoffs $J$ for each $R$
also suffices if the stated bounds and smallness hold on that set;
this is included in the formal statement. The omitted-size bound is
then needed only at those same cutoffs. For $0<p\le1$, the power-law
case is already covered by the no-envelope criterion; the broader
exponent range removes an unnecessary restriction rather than adding
a new obstruction in that range. The theorem requires these uniform
sharing bounds as hypotheses; it does not establish them for arbitrary
tape computations.

\begin{corollary}[Necessary global continuation density]\label{cor:sharing-density}
\formal{SourceCovers/SharingDensity.lean}{SourceCovers.NearFullSequence.JointProfile.seedLiminf,SourceCovers.NearFullSequence.JointProfile.necessary_sharing_density,SourceCovers.NearFullSequence.JointProfile.sharingCredit_default,SourceCovers.NearFullSequence.JointProfile.sharingCredit_eventually_pos,SourceCovers.NearFullSequence.JointProfile.sharingCost_pos,SourceCovers.NearFullSequence.JointProfile.necessary_seed_growth,SourceCovers.NearFullSequence.JointProfile.necessary_power_seed_growth,SourceCovers.NearFullSequence.JointProfile.necessary_superlinear_seed_growth,SourceCovers.NearFullSequence.JointProfile.seedLiminf_eq_top_of_zero}
Fix a near-full sequence with $c,\beta>0$ and $N\ge1$, and restrict it
to a subsequence realizing a joint profile. All dimension limits below
are along that subsequence and take values in $[0,\infty]$. For fixed $K$
with $V_K<1$ and $R\ge\max(K,1)$, define
\[
 \mu_R=\frac{(1-V_K)^2}{8(2K+1)(R+2)}-
                   \frac{\beta\tau_R}{R+2},\qquad
 g^{\rm low}_{J,R}=\liminf_{\nu\to\infty}\overline G_{J,R}^{(\nu)}.
\]
If $\mu_R>0$, then for every positive sequence $\eta_J\to0$ with
$h_{J+1}\le\eta_J$ for all sufficiently large $J$,
\begin{equation}\label{eq:sharing-density}
 \liminf_{J\to\infty}\eta_Jg^{\rm low}_{J,R}
       \ge\frac{\mu_R^2}{4A_RB_R}.
\end{equation}
The condition $\mu_R>0$ holds for all sufficiently large fixed $R$.
No bound uniform in dimension is required: it suffices to bound the
limiting hop profile. For each such $R$, without any size envelope,
the stronger conclusion is
\begin{equation}\label{eq:sharing-superlinear}
 \frac{g^{\rm low}_{J,R}}{J+1}\longrightarrow\infty.
\end{equation}
At the same fixed $R$, if $h_{J+1}\le H/(J+1)^p$ eventually,
where $H,p>0$, then
\[
 \liminf_{J\to\infty}\frac{g^{\rm low}_{J,R}}{(J+1)^p}
 \ge\frac{\mu_R^2}{4A_RB_RH}.
\]
In particular an inverse-square omitted-size bound forces positive
quadratic growth. Any common finite-dimensional size bound passes to
the limiting profile and therefore remains a sufficient hypothesis.
\end{corollary}
\begin{outline}
Fix $J$ and $\lambda>0$. Divide the finite sharing inequality by its
positive seed-count coefficient before taking the dimension lower limit.
The paid-error and width limits then give
\[
 g^{\rm low}_{J,R}\ge
 \left(\frac{\mu_{J,R}-2\lambda-A_R\eta_J/\lambda}
                  {\lambda B_R}\right)_+,
\]
where $\mu_{J,R}$ is the left side of
\cref{eq:sharing-prefix-limit}. To use only $h_{J+1}\le\eta_J$,
first replace $\eta_J$ by $\eta_J+\varepsilon$: coordinate convergence
makes this bound valid in all sufficiently large dimensions. Continuity
then removes $\varepsilon>0$. No further subsequence is needed.

For finite $g^{\rm low}_{J,R}$ and positive $\mu_{J,R}$, the shared
quadratic-tradeoff lemma gives
$\mu_{J,R}^2\le4(2+B_Rg^{\rm low}_{J,R})A_R\eta_J$.
Including nonpositive credit by its positive part yields
\[
 \eta_Jg^{\rm low}_{J,R}
 \ge\frac{(\mu_{J,R})_+^2}{4A_RB_R}-\frac{2\eta_J}{B_R}.
\]
An infinite inner lower limit satisfies this bound because $\eta_J>0$.
Now take the lower limit in $J$, using $\mu_{J,R}\to\mu_R$.
Also $\tau_R\to0$ makes $\mu_R>0$ at every sufficiently large $R$.
Extended-real scaling gives the power-growth statement.

For the stronger no-envelope conclusion, nonnegativity, monotonicity
and summability of the hop profile imply $(J+1)h_{J+1}\to0$.
The shared proof compares $(J+1)h_{J+1}$ with twice any fixed tail,
once $J$ is at least twice that tail's starting index.
Thus, for every $H>0$, the rate $H/(J+1)$ eventually bounds
$h_{J+1}$. The power-growth inequality at $p=1$ gives a lower bound
$\mu_R^2/(4A_RB_RH)$ for the normalized seed liminf. Letting $H$
decrease to zero proves \cref{eq:sharing-superlinear}.
\end{outline}

The formal density inequality also permits any $M\ge2$ with $R<M$
and independent $K,R$, and retains the positive part of the limiting
credit. At a fixed $J\ge\max(K,R)$, if $h_{J+1}=0$ and
$\mu_{J,R}>0$, its separate zero-cost case gives
$g^{\rm low}_{J,R}=\infty$.

Consequently a bound $\overline G_{J,R}\le C_R'(J+1)$, with finite
constants $C_R'$ independent of dimension, at arbitrarily large protected
scales cannot coexist with a near-full sequence. This improves the
sublinear no-envelope obstruction above.
The count concerns globally incomparable complete continuations;
it is not an information-capacity or native-gate lower bound.
No such seed-count upper bound is established for arbitrary tape computations.

\subsection{Raw metric covers instead of temporal counts}

For $\delta>0$, let $\mathcal N_{\rho,J,R}(\delta)$ be the minimum
number of closed real intervals of length at most $\delta L$ covering
$S^{\rm new}_\rho$. An empty set has covering number zero. These are
intervals in the linear raw-coordinate range, not circular arcs.

\begin{proposition}[Metric shared-support criterion]\label{prop:metric-sharing}
\formal{SourceCovers/MetricSharing.lean}{SourceCovers.metric_sharing_uniform_deficit_of_limits,SourceCovers.metric_sharing_uniform_deficit,SourceCovers.AdmissibleGeometry.prefix_metric_deficit,SourceCovers.NearFullSequence.JointProfile.prefix_metric_limit,SourceCovers.RawIntervalCover.blockUnion_card_le,SourceCovers.RawIntervalCover.two_block_factor_sharp}
In the setting of \cref{thm:uniform-sharing}, replace its bounds on
$\overline G_{J,R}$ by the following. For each $R$ in a fixed unbounded
set, there are dimension-independent numbers $\delta_{J,R}>0$
and $n_{J,R}\ge0$ with
\[
 \mathbb E_\rho\mathcal N_{\rho,J,R}(\delta_{J,R})\le n_{J,R}.
\]
Along a common unbounded set of cutoffs $J$ for that $R$, assume jointly
\begin{equation}\label{eq:metric-sharing-limits}
 \delta_{J,R}\to0,\qquad
 \delta_{J,R}n_{J,R}\to0,\qquad
 \eta_J/\delta_{J,R}\to0.
\end{equation}
Then there are $\eps>0,d_0$ such that $a_0\le\eps L$ and $d\ge d_0$
imply $\Delta\ge\eps L$ for this class.
More generally, it suffices that for each fixed $R$, every $\varepsilon>0$
and every $J_0$ admit a common allowed cutoff $J\ge J_0$ at which
$\delta_{J,R}$, $\delta_{J,R}n_{J,R}$ and $\eta_J/\delta_{J,R}$ are
all smaller than $\varepsilon$. The geometric bounds need only hold at
the allowed cutoffs; these sets remain independent of dimension.
\end{proposition}
\begin{outline}
Singleton intervals cover any finite raw set, so the set of achievable
cover cardinalities has a least element, attained by a cover. This uses
natural-number minimization; real endpoints and overlapping intervals
are permitted.

Put $w=\lceil\delta_{J,R}L\rceil>0$. A set contained in an interval of
length at most $\delta_{J,R}L\le w$ has grid labels differing by at most
one. It therefore meets at most two blocks, each containing at most $w$
raw coordinates. The union of all covered blocks has cardinality at most
\[
 \min\{L,\,2w\mathcal N_{\rho,J,R}(\delta_{J,R})\}
\]
at every shift, without a disjointness assumption on the cover. The factor
two is sharp for this finite estimate: width-one blocks and the interval
$[0,1]$ covering the two raw coordinates $\{0,1\}$ attain it.

The generalized raw-payment theorem accepts any such shift-uniform bound.
It retains the full two-exception residual bill and the older bank's
movement bound; the original per-coordinate payment is a direct
specialization. Average the new payment over the same phases and use the
existing rounded-width limit, with all cutoffs and $\delta_{J,R}$ fixed
before dimension. The right side of \cref{eq:sharing-prefix-limit} becomes
\[
 2\delta_{J,R}+2\delta_{J,R}n_{J,R}
              +A_R\eta_J/\delta_{J,R}.
\]
Simultaneous smallness at arbitrarily late cutoffs makes the limiting
credit at most zero after subtracting $\beta\tau_R/(R+2)$. The same
vanishing-tail contradiction used in \cref{thm:uniform-sharing} completes
the proof. The three joint limits in \cref{eq:metric-sharing-limits}
imply this smallness condition on their common unbounded set.
\end{outline}

The improvement from three blocks to two sharpens the finite payment
and its limiting coefficient; it does not change the vanishing conditions
or establish metric covers for an unrestricted native class.

\begin{corollary}[Sublinear raw covering exponents]\label{cor:covering-exponents}
\notformalized
The metric criterion holds in either of the following classes, with
constants uniform in dimension and cutoff and permitted to depend on
each fixed $R$ in an unbounded set.
\begin{enumerate}[label=(\alph*)]
\item Without a size envelope,
$\mathbb E_\rho\mathcal N_{\rho,J,R}(\delta)
 \le H_R\delta^{-s_R}$ for all $J\ge R$ and
$0<\delta<\delta_{0,R}$, where $H_R,\delta_{0,R}>0$ and $0\le s_R<1$.
\item Under $|G_i|\le HL/i^p$ with fixed $p>1$,
$\mathbb E_\rho\mathcal N_{\rho,J,R}(\delta)
 \le H_RJ^{u_R}\delta^{-s_R}$ in the same range, where
$0\le s_R<1$ and $0\le u_R<p(1-s_R)$.
\end{enumerate}
\end{corollary}
\begin{informalproof}
In (a) take $\eta_J=c/(J+1)$ and
$\delta_{J,R}=(J+1)^{-1/2}$. In (b) choose a fixed exponent
$u_R/(1-s_R)<v_R<p$ and set
$\delta_{J,R}=(J+1)^{-v_R}$, with
$\eta_J=\min\{c/(J+1),H/(J+1)^p\}$.
The products $\delta n$ are respectively bounded by constants times
$(J+1)^{-(1-s_R)/2}$ and $J^{u_R-v_R(1-s_R)}$; the residual ratios
$\eta_J/\delta$ also tend to zero. Eventually $\delta<\delta_{0,R}$.
Apply \cref{prop:metric-sharing}.
\end{informalproof}

These conditions concern raw-coordinate support shared by all outputs.
They can be more economical than counting separate temporal occurrences.
Supplying them for a still-open native class remains a geometric problem;
the hypotheses are not consequences of a low entropy estimate alone.
